\section*{Chapitre \ref{ch:iv1:expectation-and-indicators} --- Espérance et variables indicatrices}
\markboth{\protect\omTocMark{}{Réponses~: Espérance et variables indicatrices}}{\protect\omTocMark{}{Réponses~: Espérance et variables indicatrices}}

\enlargethispage{-4\baselineskip}%
\begin{solution}{iq:iv1:expectation-and-indicators:1}
L'espérance vaut $161/36 \approx 4.47$.

On utilise la formule de la somme des queues\index{formule de la somme des queues}~: pour une variable prenant les valeurs $1, 2, \dots$, l'espérance est égale à
$\sum_{k \ge 1} P(M \ge k)$, car chaque valeur de $M$ est comptée une fois pour chaque $k$ jusqu'à $M$. Le dé le plus élevé vaut au
moins $k$, sauf si les deux dés sont tous deux inférieurs à $k$, donc $P(M \ge k) = 1 - \bigl(\tfrac{k-1}{6}\bigr)^2$. En sommant sur
$k = 1, \dots, 6$,
\[
E[M] = 6 - \frac{0 + 1 + 4 + 9 + 16 + 25}{36} = 6 - \frac{55}{36} = \frac{161}{36}.
\]
\iqfollow{Et avec trois dés~?}{Les cubes remplacent les carrés~: $6 - 225/216 = 119/24 \approx 4.96$.}
\iqfollow{En moyenne, quelle est la distance entre les deux dés~?}{L'écart est le dé le plus élevé moins le plus faible, et
le dé le plus élevé plus le plus faible est la somme des dés, qui vaut en moyenne 7. Le dé le plus faible a donc pour espérance $7 - 161/36$, et
l'écart vaut en moyenne $2 \cdot 161/36 - 7 = 35/18 \approx 1.94$.}
\end{solution}

\begin{solution}{iq:iv1:expectation-and-indicators:2}
L'espérance du total vaut 21.

Le nombre de lancers $N$ suit une loi géométrique d'espérance 6, et chaque lancer a pour espérance $7/2$. Le fait de s'arrêter après un lancer
ne dépend que des lancers précédents, donc l'identité de Wald\index{identité de Wald} s'applique~: $E[\text{total}] = E[N] \cdot
7/2 = 21$. On peut le vérifier avec l'analyse au premier pas\index{analyse au premier pas}. Un lancer qui n'est pas un six a pour espérance 3, donc
\[
S = \tfrac16 \cdot 6 + \tfrac56 \,(3 + S),
\]
ce qui donne $S = 21$.

\iqwrong{«~Cinq lancers précédents de moyenne 3.5, plus un six, font 23.5.~» Les lancers précédents sont connus pour ne pas être des six, donc
ils ont pour moyenne 3, et non 3.5. Cinq d'entre eux plus le six donnent $5 \cdot 3 + 6 = 21$.}
\iqfollow{Combien de 1 obtenez-vous en moyenne avant le 6~?}{Exactement un. On ignore chaque lancer sauf les 1 et les 6~:
chacun de ces deux résultats a autant de chances de sortir ensuite, donc le nombre de 1 avant le premier 6 suit une loi géométrique d'espérance 1.}
\end{solution}

\begin{solution}{iq:iv1:expectation-and-indicators:3}
En moyenne, il y a exactement 3 paires de ce type.

On utilise la linéarité de l'espérance avec des variables indicatrices\index{linéarité de l'espérance}\index{variables indicatrices}. Il y a 51
paires voisines. Quelle que soit la première carte d'une paire, la seconde est l'une des 51 autres cartes, et 3 d'entre elles partagent
sa valeur, donc chaque paire est assortie avec une probabilité $3/51$. Le total vaut $51 \cdot 3/51 = 3$. Les événements se recouvrent et sont dépendants,
et cela ne change rien~: la linéarité n'exige aucune indépendance.

\iqfollow{Quelle est la variance de ce nombre~?}{On somme sur les couples de paires. Il existe 50 couples de paires voisines qui se chevauchent, et chacun
s'assortit avec la probabilité $\frac{3}{51}\cdot\frac{2}{50} = \frac{1}{425}$. Il y a 1\,225 couples de paires disjointes, chacun avec la
probabilité $\frac{3}{51}\cdot\frac{146}{2450}$, où $\frac{146}{2450} = \frac{2}{50}\cdot\frac{1}{49} +
\frac{48}{50}\cdot\frac{3}{49}$ se décompose selon que la première carte de la seconde paire a la valeur qui vient d'être utilisée. Puis $\mathrm{Var} = 3 + 2
\cdot \frac{50}{425} + 2\cdot 1225\cdot\frac{3}{51}\cdot\frac{146}{2450} - 9 = \frac{48}{17} \approx 2.82$.}
\iqfollow{Quelle est la probabilité qu'un mélange ne contienne aucune paire de ce type~?}{Environ 4.5~\%, par simulation. Une loi de Poisson d'espérance 3 prédit
$e^{-3} \approx 5~\%$, ce qui est proche car les nombreuses variables indicatrices sont rares et presque indépendantes.}
\end{solution}

\enlargethispage{-3\baselineskip}%
\begin{solution}{iq:iv1:expectation-and-indicators:4}
L'espérance vaut 1 et la variance vaut 1, pour un jeu de cartes de toute taille supérieure à une carte.

On utilise des variables indicatrices\index{variables indicatrices}~: la position $i$ est un point fixe avec probabilité $1/52$, donc l'espérance vaut $52 \cdot 1/52 = 1$.
Pour la variance, on écrit $X^2$ comme la somme des variables indicatrices plus la somme sur les couples ordonnés de positions de $I_iI_j$. Les positions
$i$ et $j$ sont toutes deux des points fixes avec probabilité $\frac{1}{52 \cdot 51}$, et il y a $52 \cdot 51$ couples ordonnés, donc $E[X^2] = 1 + 1 =
2$ et $\mathrm{Var}(X) = 2 - 1 = 1$.

\iqwrong{Traiter les positions comme indépendantes donne une variance de $52 \cdot \frac{1}{52}\cdot\frac{51}{52} = 51/52$. La petite
covariance positive entre les positions fournit le reste.}
\iqfollow{Quelle est la probabilité de n'avoir aucun point fixe~?}{Très proche de $1/e \approx 0.368$, le nombre de points fixes étant presque de loi de Poisson
de moyenne 1. C'est la probabilité d'un dérangement\index{dérangement}.}
\end{solution}

\begin{solution}{iq:iv1:expectation-and-indicators:5}
L'espérance est de $14.7$ lancers.

C'est le problème du collectionneur de coupons\index{collectionneur de coupons}. On découpe l'attente en étapes. Lorsqu'on a vu $k$ faces, chaque lancer montre une
nouvelle face avec probabilité $(6-k)/6$, si bien que l'étape dure $6/(6-k)$ lancers en moyenne. En additionnant les étapes,
\[
6\Bigl(\tfrac16 + \tfrac15 + \tfrac14 + \tfrac13 + \tfrac12 + 1\Bigr) = 14.7.
\]
\iqfollow{Quelle est la variance~?}{Les étapes sont des lois géométriques indépendantes, chacune de variance $(1-p)/p^2$. Leur somme vaut environ
$38.99$, soit un écart-type d'environ $6.2$ lancers. La queue est longue car la dernière face prend à elle seule 6 lancers en moyenne.}
\iqfollow{Quelle est l'espérance de la somme des lancers~?}{L'identité de Wald\index{identité de Wald} donne $14.7 \cdot 3.5 = 51.45$.}
\iqwhy{Le coût vaut $n \ln n$, et non $n$. Pour voir chacun des 1\,000 instruments au moins une fois dans des échantillons de données tirés uniformément au hasard,
il faut environ 7\,500 tirages, et non 1\,000. C'est ainsi qu'on dimensionne un échantillon devant toucher chaque casier.}
\end{solution}

\begin{solution}{iq:iv1:expectation-and-indicators:6}
L'espérance vaut 11 et la variance vaut 5.

Les séries se comptent plus facilement à travers les changements. Le nombre de séries vaut un de plus que le nombre de lancers voisins qui
diffèrent. Chacune des 20 paires voisines diffère avec probabilité $1/2$, donc grâce aux variables indicatrices\index{variables indicatrices}, l'espérance vaut
$1 + 20/2 = 11$. Les changements sont indépendants, et pas seulement deux à deux~: sachant un lancer quelconque, le suivant diffère avec probabilité $1/2$
quel que soit le passé. Le nombre de changements suit donc une loi binomiale $(20, 1/2)$, de variance $20/4 = 5$.

\iqfollow{Sur 100 lancers, combien de séries comportant exactement trois piles d'affilée attendez-vous~?}{Une série commençant au milieu exige le motif
FPPPF, avec probabilité $1/32$, et il y a 96 positions de départ possibles. Une série à l'une ou l'autre extrémité n'exige que PPPF ou FPPP, avec
probabilité $1/16$ chacune. Le total vaut $96/32 + 2/16 = 3.125$.}
\iqwhy{C'est le test des séries\index{test des séries} pour repérer les suites de rendements de même signe. Pour une année de 252 signes quotidiens, le nombre de séries a pour espérance
$126.5$ et un écart-type d'environ $7.9$, si bien qu'un nombre inférieur à environ 111 constitue un signal à deux écarts-types de suites.}
\end{solution}

\begin{solution}{iq:iv1:expectation-and-indicators:7}
Le deuxième as apparaît en moyenne à la position $106/5 = 21.2$.

On utilise un argument de symétrie\index{arguments de symétrie}. Les 4 as découpent les 48 autres cartes en 5 intervalles, comme dans la \cref{fig:iv1:ace-gaps}.
Chacune de ces 48 cartes a autant de chances de tomber dans un intervalle que dans un autre, car elle peut se placer à chacune des 5 places
possibles par rapport aux 4 as, avec la même probabilité. Chaque intervalle contient donc $48/5 = 9.6$ cartes en moyenne. Le second as vient après deux intervalles et deux as~:
\[
2 \cdot \frac{48}{5} + 2 = \frac{106}{5}.
\]
\begin{omfigure}
\begin{tikzpicture}[x=0.9cm,y=0.9cm]
\foreach \a/\b/\l in {0/2/G_1,2.4/4.4/G_2,4.8/6.8/G_3,7.2/9.2/G_4,9.6/11.6/G_5}{
  \fill[omDef!15] (\a,0) rectangle (\b,0.5);
  \draw[omDef] (\a,0) rectangle (\b,0.5);
  \node at ({(\a+\b)/2},0.25) {\small $\l$};
}
\foreach \x/\n in {2.2/1,4.6/2,7/3,9.4/4}{
  \fill[omThm] (\x-0.17,0) rectangle (\x+0.17,0.5);
  \node[above,font=\footnotesize] at (\x,0.5) {as \n};
}
\draw[decorate,decoration={brace,mirror,amplitude=5pt}] (0,-0.08) -- (4.75,-0.08)
  node[midway,below=5pt,font=\footnotesize] {deux intervalles et deux as~: $2\cdot 9.6 + 2 = 21.2$};
\end{tikzpicture}

\omcaption{Les quatre as découpent le jeu en cinq intervalles d'espérance égale. Le deuxième as se place après deux intervalles et les deux as
eux-mêmes.}\label{fig:iv1:ace-gaps}
\end{omfigure}

\iqwrong{$2 \cdot 52/5 = 20.8$ oublie que les as occupent aussi des positions.}
\iqfollow{En moyenne, où se trouve le dernier as~?}{Le même calcul donne le $k$-e as à la position $k \cdot 53/5$, donc le quatrième
se trouve en $42.4$.}
\end{solution}

\begin{solution}{iq:iv1:expectation-and-indicators:8}
L'espérance vaut exactement $25.5025$.

On utilise la formule de la somme des queues\index{formule de la somme des queues}. Le plus petit des trois dés vaut au moins $k$ lorsque les trois dés valent au moins $k$, ce qui a pour
probabilité $\bigl(\tfrac{101-k}{100}\bigr)^3$. En sommant de $k = 1$ à $100$ et en posant $j = 101 - k$,
\[
E[\min] = \frac{1}{100^3} \sum_{j=1}^{100} j^3 = \frac{1}{100^3}\Bigl(\frac{100 \cdot 101}{2}\Bigr)^2 = \frac{5050^2}{10^6} = 25.5025,
\]
en utilisant le fait que la somme des $n$ premiers cubes est le carré de la somme des $n$ premiers entiers.

\iqwrong{En appliquant la règle continue, le plus petit de trois uniformes se situe au quart de l'intervalle et donne 25. C'est proche, mais elle
rate le surplus de $0.5025$ dû aux faces discrètes. Un recruteur qui demande la valeur exacte attend la somme.}
\iqfollow{Quelle est l'espérance du plus grand des trois~?}{Par la symétrie qui envoie une face $f$ sur $101 - f$, elle vaut $101 - 25.5025 = 75.4975$.}
\end{solution}

\begin{solution}{iq:iv1:expectation-and-indicators:9}
Pour le dé équilibré, l'espérance du nombre de lancers vaut $1223/324 \approx 3.77$. Pour la table de hachage, environ 41 clés.

On utilise la formule de la somme des queues\index{formule de la somme des queues}, $E[T] = \sum_{k \ge 0} P(T > k)$. L'attente dépasse $k$ lorsque les $k$ premiers lancers sont tous différents, ce qui a pour
probabilité $\frac{6}{6}\cdot\frac{5}{6}\cdots\frac{6-k+1}{6}$. Les termes pour $k = 0, 1, \dots, 6$ sont
\[
1,\quad 1,\quad \tfrac56,\quad \tfrac59,\quad \tfrac{5}{18},\quad \tfrac{5}{54},\quad \tfrac{5}{324},
\]
et leur somme vaut $1223/324$. Remarquons que l'attente vaut au moins 2 et au plus 7.

Pour $n$ compartiments, la même somme est le problème des anniversaires\index{problème des anniversaires}. Sa valeur pour $n$ grand vaut environ $\sqrt{\pi n/2} + 2/3$, ce qui, pour
$n = 1024$ donne $40.8$~; la somme exacte vaut elle aussi $40.8$. Ce décompte inclut la clé qui provoque la collision.

\iqwhy{Les collisions arrivent à la racine carrée de la taille de la table, et non à la taille elle-même. Un identifiant d'ordre de 32~bits choisi au hasard a
sa première collision attendue après environ 82\,000 ordres, et non 4 milliards.}
\end{solution}

\enlargethispage{-3\baselineskip}%
\begin{solution}{iq:iv1:expectation-and-indicators:10}
En moyenne, on tire $e \approx 2.718$ nombres.

Soit $N$ le nombre de tirages. Alors $N > n$ exactement lorsque les $n$ premiers nombres ont une somme au plus égale à 1, et cela a pour probabilité $1/n!$. Pour le voir
il suffit de passer aux sommes partielles $s_1 < s_2 < \dots < s_n$. L'application a pour déterminant un, donc la probabilité est le volume du
domaine $0 < s_1 < \dots < s_n < 1$, et ce domaine est l'un des $n!$ ordres égaux de $n$ points dans le cube unité. On applique la
formule de la somme des queues\index{formule de la somme des queues} donne
\[
E[N] = \sum_{n \ge 0} P(N > n) = \sum_{n \ge 0} \frac{1}{n!} = e.
\]
\iqwrong{«~Chaque nombre vaut en moyenne un demi, donc deux nombres.~» Cela ignore le dépassement~: le total au moment de l'arrêt est supérieur à 1, et non égal à 1.}
\iqfollow{Quel est le total moyen au moment où vous arrêtez~?}{Par l'identité de Wald\index{identité de Wald}, $E[N] \cdot \tfrac12 = e/2 \approx 1.359$.}
\iqfollow{Quelle est la variance de $N$~?}{$E[N^2] = \sum_{n \ge 0} (2n+1)P(N>n) = 2e + e = 3e$, donc la variance vaut $3e - e^2 \approx 0.766$.}
\end{solution}

\begin{solution}{iq:iv1:expectation-and-indicators:11}
L'espérance du nombre de tirages supplémentaires est infinie, et pourtant un seul tirage supplémentaire suffit une fois sur deux.

Soit $T$ le nombre de tirages supplémentaires. $T > k$ signifie que le premier nombre est le plus grand des $k+1$ premiers, et par symétrie
cet événement a pour probabilité $1/(k+1)$. Ainsi $P(T = 1) = 1/2$, et $P(T \le 3) = 3/4$. La formule de la somme des queues\index{formule de la somme des queues}
donne la série harmonique,
\[
E[T] = \sum_{k \ge 0} \frac{1}{k+1} = \infty.
\]
Le rare premier nombre, très grand, produit des attentes très longues, et ces attentes sont assez longues pour compenser leur rareté.

\iqfollow{Et si vous abandonnez au plus tard après 1\,000 tirages~?}{La moyenne vaut alors la somme harmonique $H_{1000} \approx 7.49$. Si l'on modifie la
limite, la réponse change avec elle, ce qui est la signature d'une moyenne divergente.}
\iqwhy{L'attente moyenne mesurée sur un échantillon ne se stabilise jamais~: on quadruple l'échantillon et elle continue de croître. Si l'on mesure le «~temps pour retrouver le plus haut précédent~» sur un
backtest sans dérive, la moyenne empirique ne renseigne que sur la durée du backtest. Il vaut mieux donner la médiane ou un quantile.}
\end{solution}

\begin{solution}{iq:iv1:expectation-and-indicators:12}
La moyenne est le nombre harmonique $H_{100} \approx 5.19$, et la variance vaut environ $3.55$.

On utilise des variables indicatrices\index{variables indicatrices}. Le $k$-e nombre est un record lorsqu'il est le plus grand des $k$ premiers, avec probabilité $1/k$, donc
$E[R] = \sum_{k=1}^{100} 1/k \approx 5.187$. La forme pour $n$ grand est $\ln n + 0.577$, soit ici environ $5.18$.

Pour la variance, il faut un résultat plus fort~: les indicatrices sont indépendantes. Savoir comment les $k-1$ premiers nombres sont ordonnés entre
eux ne dit rien de la place du $k$-e parmi les $k$ premiers, et les événements de record ne dépendent que de ces rangs relatifs. Donc la
variance est la somme des variances de Bernoulli individuelles~:
\[
\mathrm{Var}(R) = \sum_{k=1}^{100} \frac1k\Bigl(1 - \frac1k\Bigr) \approx 5.187 - 1.635 = 3.55.
\]
\iqwhy{Les records dans du bruit pur sont fréquents au début et rares ensuite~: une année de 252 résultats quotidiens indépendants compte environ $6.1$ jours qui établissent un nouveau plus haut.
Une stratégie qui établit six nouveaux plus hauts en un an n'a donc encore rien montré.}
\end{solution}

\enlargethispage{-3\baselineskip}%
\begin{solution}{iq:iv1:expectation-and-indicators:13}
La covariance vaut $-5/3$ et la corrélation vaut $-1/5$.

Posons $A = \sum_i I_i$ et $B = \sum_i J_i$, où $I_i$ indique que le lancer $i$ affiche un, et $J_i$ que le lancer $i$ affiche deux. Des lancers différents sont
indépendants, donc seul un même lancer contribue à la covariance. Un seul lancer ne peut pas être les deux à la fois, donc $I_iJ_i = 0$, et
\[
\mathrm{Cov}(I_i, J_i) = 0 - \tfrac16 \cdot \tfrac16 = -\tfrac{1}{36}.
\]
Sur 60 lancers, cela donne $-60/36 = -5/3$. Chacun des deux comptes a pour variance $60 \cdot \tfrac16 \cdot \tfrac56 = 25/3$, donc la corrélation vaut
$(-5/3)/(25/3) = -1/5$, quel que soit le nombre de lancers.

\iqfollow{Vérifiez-le autrement.}{$A + B$ compte les lancers qui affichent un ou deux, nombre qui suit une loi binomiale $(60, 1/3)$ de variance $40/3$. Et en effet
$25/3 + 25/3 - 2\cdot 5/3 = 40/3$.}
\iqwhy{Des comptes répartis dans des compartiments de total fixé sont corrélés négativement par construction. Si l'on répartit des ordres sur 6 plateformes équivalentes,
les parts de chaque plateforme ont une corrélation de $-1/5$ sans aucun lien économique entre les plateformes, donc il ne faut pas lire une corrélation négative de cette taille comme
une concurrence.}
\end{solution}

\begin{solution}{iq:iv1:expectation-and-indicators:14}
L'espérance du nombre de lancers vaut $3/2$. Ce n'est pas 3.

Considérons plutôt d'arrêter au premier lancer dans $\{1, 3, 5, 6\}$. Ce temps d'arrêt $T$ suit une loi géométrique de probabilité de succès $4/6$, donc $E[T] = 3/2$.
Votre événement «~tous les lancers ont été pairs et le dernier est un six~» est exactement l'événement où ce premier lancer d'arrêt est un six. La valeur du lancer d'arrêt est
indépendante de l'instant où il survient, puisque chacune des quatre faces est équiprobable à tout instant. Conditionner sur le fait qu'il s'agit d'un six ne modifie donc pas la
loi de $T$\index{conditionnement sur la valeur d'arrêt}.

Un calcul direct le confirme. La probabilité de s'arrêter au lancer $k$ avec tous les lancers pairs vaut $(1/3)^{k-1}(1/6)$. En sommant sur $k$, on obtient $1/4$, probabilité de la condition, et en sommant $k$
fois ce terme, on obtient $3/8$, de sorte que l'espérance conditionnelle vaut $(3/8)/(1/4) = 3/2$.

\iqwrong{«~Sachant que seuls des pairs sont sortis, le dé n'a plus que trois faces, donc l'attente d'un six suit une loi géométrique de $p = 1/3$, d'espérance 3.~» C'est lancer un dé à trois faces, ce qui
est une autre expérience. Dans la nôtre, une longue suite doit aussi éviter tous les nombres impairs, si bien que la condition rend les longues suites rares.}
\iqfollow{Et si vous saviez seulement qu'aucun lancer n'a été effectué~?}{On s'arrête sur $\{1, 6\}$, de probabilité de succès $1/3$, donc l'espérance vaut 3.}
\iqwhy{Conditionner sur la façon dont un chemin se termine modifie la durée qu'elle semble avoir eue. Prenons une marche aléatoire équilibrée avec un stop à 10 en dessous et un objectif à 5 au-dessus~: tous les trades durent en moyenne 50 pas, mais les gagnants durent $41.7$ et les perdants $66.7$. Un rapport de backtest sur le «~temps jusqu'à l'objectif~» calculé sur les seuls gagnants sous-estime la durée de détention.}
\end{solution}

\begin{solution}{iq:iv1:expectation-and-indicators:15}
Le nombre espéré de boucles vaut $\sum_{k=1}^{100} \frac{1}{2k-1} \approx 3.28$.

Conditionnons sur la première étape\index{analyse au premier pas}. Prenons $n$ nouilles, soit $2n$ bouts libres. Choisissons un bout~; son partenaire est équiprobable parmi les
$2n - 1$ autres bouts. Exactement un de ceux-ci est l'autre bout de la même nouille, et une boucle se referme alors. Dans tous les autres cas, deux nouilles fusionnent en
une nouille plus longue, qui se comporte exactement comme une seule nouille. Dans les deux cas, il reste $n - 1$ nouilles, donc
\[
L_n = L_{n-1} + \text{Bernoulli}\Bigl(\frac{1}{2n-1}\Bigr), \qquad E[L_n] = \sum_{k=1}^{n} \frac{1}{2k-1}.
\]
Pour $n$ grand, cela vaut environ $\tfrac12 \ln n + 0.98$, et pour $n = 100$ cela donne $3.28$. Choisir deux bouts au hasard revient à choisir un bout
puis un partenaire au hasard, de sorte que l'ordre adopté ne fait rien perdre.

\iqwrong{Deviner un nombre proportionnel à $n$. La plupart des liaisons fusionnent des nouilles au lieu de les refermer, et le nombre de boucles croît comme le logarithme de $n$.}
\iqfollow{Quelle probabilité accorderiez-vous à ce que tout se termine en une seule grande boucle~?}{Aucune liaison ne referme une boucle avant la toute dernière~: $\prod_{k=2}^{100} \frac{2k-2}{2k-1} \approx 0.089$.}
\end{solution}

\begin{solution}{iq:iv1:expectation-and-indicators:16}
La dernière fourmi tombe au bout de $100/101 \approx 0.990$ minute en moyenne, et il y a $1237.5$ collisions en moyenne.

Les fourmis étant identiques, une collision où les deux font demi-tour ressemble exactement à deux fourmis qui se traversent, comme dans la
\cref{fig:iv1:ants}. Suivons ces «~fantômes~»\index{réétiquetage d'objets identiques}. Chaque fantôme marche droit vers le bout auquel il fait face et sort. Son temps de marche est la distance à ce bout. Pour
une position uniforme et une direction équilibrée, cette distance est uniforme sur $[0, 1]$, indépendamment pour chaque fourmi. La dernière chute est le maximum de 100
uniformes indépendantes, dont l'espérance vaut $n/(n+1) = 100/101$.

Chaque collision réelle est une traversée de deux fantômes. Deux fantômes se croisent exactement lorsque celui de gauche regarde à droite et celui de droite regarde à gauche, avec probabilité
$1/4$, et ils se rencontrent alors toujours avant que l'un ou l'autre ne tombe. Le nombre espéré vaut donc $\binom{100}{2}/4 = 1237.5$.

\begin{omfigure}
\begin{tikzpicture}[x=7cm,y=4.8cm]
\draw[->] (0,0) -- (1.08,0) node[right,font=\footnotesize] {position};
\draw[->] (0,0) -- (0,0.92) node[above,font=\footnotesize] {temps};
\draw[thick] (1,0) -- (1,0.85);
\node[below,font=\footnotesize] at (0,0) {0};
\node[below,font=\footnotesize] at (1,0) {1};
\draw[omDef,very thick] (0.25,0) -- (0.475,0.225) -- (0,0.7);
\draw[omThm,very thick] (0.7,0) -- (0.475,0.225) -- (1,0.75);
\node[right,font=\footnotesize,xshift=6pt] at (0.475,0.225) {collision};
\node[left,font=\footnotesize,omDef] at (0,0.7) {la fourmi 1 tombe à $0.70$};
\node[right,font=\footnotesize,omThm] at (1,0.75) {la fourmi 2 tombe à $0.75$};
\node[below,font=\footnotesize,omDef] at (0.25,0) {fourmi 1};
\node[below,font=\footnotesize,omThm] at (0.7,0) {fourmi 2};
\end{tikzpicture}

\omcaption{Deux fourmis sur le bâton, le temps croissant vers le haut. Chaque fourmi fait demi-tour à la collision et reste de son côté. Si l'on oublie les couleurs, la même figure
montre deux droites qui se traversent, les fantômes. Les positions et les instants de chute sont les mêmes dans les deux lectures.}\label{fig:iv1:ants}
\end{omfigure}

\iqfollow{Diriez-vous que le temps de chute d'un fantôme appartient à la fourmi dont il est issu~?}{Non, et la question n'en avait pas besoin. Les temps de chute sont justes en tant qu'ensemble, pas fourmi par fourmi~:
les vraies fourmis conservent leur ordre de gauche à droite, donc celles qui sortent à gauche sortent dans l'ordre de leur départ, et celles qui sortent à droite dans l'ordre inverse.}
\end{solution}

\begin{solution}{iq:iv1:expectation-and-indicators:17}
En moyenne, on retourne $689/45 \approx 15.31$ cartes.

Soit $A$ la position du premier as et $K$ celle du premier roi. On s'arrête en $\max(A, K)$. Utilisons l'identité max-min\index{inclusion-exclusion pour les maxima}~:
\[
\max(A, K) = A + K - \min(A, K).
\]
Le premier de $m$ cartes spéciales dans un jeu de 52 cartes se trouve en moyenne en $53/(m+1)$, par symétrie des intervalles~: $m$ cartes spéciales découpent les autres en $m+1$ intervalles de
même espérance, et la première carte spéciale suit l'un d'eux. Avec 4 as, cela donne $E[A] = E[K] = 53/5$. Avec les 8 as et rois réunis, on obtient
$E[\min] = 53/9$. Donc
\[
E[\max(A, K)] = \frac{2 \cdot 53}{5} - \frac{53}{9} = \frac{689}{45}.
\]
\iqwrong{Additionner les deux attentes, $53/5 + 53/5$. Elles courent en même temps, et elles se recouvrent exactement du temps d'attente de la première carte spéciale.}
\iqfollow{Un dé à trois faces biaisé montre ses faces avec les probabilités $\tfrac12, \tfrac13, \tfrac16$. Combien de lancers vous faut-il pour avoir vu les trois faces~?}{Inclusion-exclusion
sur les sous-ensembles, avec une attente de $1/p$ pour tout ensemble de probabilité totale $p$~: $2 + 3 + 6 - \frac65 - \frac32 - 2 + 1 = 7.3$.}
\iqwhy{C'est ainsi qu'on calcule le temps nécessaire pour que les deux jambes d'un ordre apparié soient exécutées. Si chaque jambe prend 10 secondes en moyenne, de façon indépendante et sans mémoire, alors les deux sont terminées au bout de
$10 + 10 - 5 = 15$ secondes, ni 10 ni 20.}
\end{solution}

\begin{solution}{iq:iv1:expectation-and-indicators:18}
Le nombre espéré de chaussettes vaut $256/63 \approx 4.06$.

Utilisons la formule de la queue de distribution\index{formule de la queue de distribution}. Ne pas avoir de paire après $k$ chaussettes signifie qu'elles proviennent de $k$ paires différentes. Il y a $\binom{5}{k}$ façons de choisir les paires et $2^k$ façons
de prendre une chaussette de chacune, sur $\binom{10}{k}$ tirages possibles~:
\[
P(T > k) = \frac{2^k \binom{5}{k}}{\binom{10}{k}}, \qquad k = 0, 1, \dots, 5.
\]
Ces valeurs sont $1, 1, \tfrac89, \tfrac23, \tfrac{8}{21}, \tfrac{8}{63}$, et leur somme vaut $256/63$. La réponse se situe entre 2 et 6, comme il se doit.

Pour $n$ paires, la somme admet la forme close $4^n / \binom{2n}{n}$, qui croît comme $\sqrt{\pi n}$. Pour $n = 5$, l'approximation donne $3.96$,
déjà proche.

\iqwrong{Utiliser la formule des anniversaires avec $5$ types de chaussettes équiprobables donne environ $3.51$, ce qui est trop bas. Après la première chaussette, seule 1 des 9 chaussettes restantes lui correspond, soit une chance de $1/9$ et non $1/5$.}
\iqfollow{Et si vous aviez 10 paires~?}{$4^{10}/\binom{20}{10} \approx 5.68$. Doubler le nombre de chaussettes multiplie l'attente par environ $\sqrt{2}$.}
\end{solution}

\begin{solution}{iq:iv1:expectation-and-indicators:19}
En moyenne, il reste $62/35 \approx 1.77$ boules.

Utilisons une variable indicatrice\index{variables indicatrices} pour chaque boule, et demandons-nous quand elle se trouve encore dans l'urne à la fin. Une boule rouge reste exactement lorsque toutes les boules bleues sont tirées avant elle. Ne regardons que
l'ordre relatif de cette boule rouge et des 4 boules bleues~: les 5 ordres possibles de «~qui vient en dernier~» sont équiprobables, donc cela arrive avec probabilité $1/5$. De même, une boule bleue reste lorsque les 6 boules rouges passent toutes avant elle, avec probabilité
$1/7$. En sommant sur les boules,
\[
E[\text{restantes}] = \frac{6}{5} + \frac{4}{7} = \frac{62}{35}.
\]
Plus généralement, avec $r$ boules rouges et $b$ boules bleues, cela vaut $\frac{r}{b+1} + \frac{b}{r+1}$.

\iqwrong{Modéliser le tirage avec des probabilités fixes et une attente géométrique. La composition de l'urne change à chaque tirage, et les variables indicatrices évitent de devoir la suivre.}
\iqfollow{Quelle est la probabilité que les boules restantes soient rouges~?}{Elles sont rouges exactement lorsque la toute dernière boule d'un tirage complet est rouge, et par symétrie cette probabilité vaut $6/10 = 3/5$.}
\iqfollow{Combien de boules tirez-vous en moyenne~?}{$10 - 62/35 = 288/35 \approx 8.23$.}
\end{solution}

\begin{solution}{iq:iv1:expectation-and-indicators:20}
La moyenne vaut environ $36.60$ et la variance environ $9.74$, donc l'écart type vaut environ $3.1$.

Utilisons des variables indicatrices\index{variance d'un comptage}. La case $i$ est vide avec probabilité $q_1 = (0.99)^{100}$, donc la moyenne vaut $100\,q_1 \approx 36.60$, soit environ $100/e$. Les cases $i$ et $j$ sont toutes deux vides avec
probabilité $q_2 = (0.98)^{100}$, puisque chaque boule doit éviter les deux. En développant le carré de la somme,
\[
\mathrm{Var}(X) = 100\,q_1 + 100 \cdot 99\, q_2 - (100\,q_1)^2 \approx 9.74.
\]
Pour $n$ grand, cela tend vers $n(e^{-1} - 2e^{-2})$, soit ici $9.72$.

\iqwrong{Traiter le comptage comme une loi binomiale, $100\,q_1(1 - q_1) \approx 23.2$. Cela surestime la variance réelle de plus du double. Les cases sont négativement corrélées, avec
$\mathrm{Cov} = q_2 - q_1^2 \approx -0.00136$. Si la case 1 est vide, ses boules sont allées ailleurs, donc les autres cases ont davantage de chances d'être remplies. C'est petit par paire, mais il y a 9\,900 paires ordonnées.}
\iqwhy{Cela donne un contrôle rapide d'une fonction de hachage. Avec 100 clés et 100 cases, 50 cases vides se situent à plus de 4 écarts types au-dessus de la moyenne. La variance binomiale
l'aurait ramené à moins de 3 écarts types, et on aurait pu accepter la fonction de hachage.}
\end{solution}

\begin{solution}{iq:iv1:expectation-and-indicators:21}
Pour 10 nombres, le nombre moyen de comparaisons vaut environ $24.44$. En général, il vaut $2(n+1)H_n - 4n$, où $H_n = \sum_{k \le n} 1/k$.

Utilisons une variable indicatrice par paire\index{variables indicatrices}. Numérotons les nombres selon leur rang dans l'ordre croissant, $1, \dots, n$, et fixons $i < j$. Deux nombres sont comparés au plus une fois, car une comparaison fait de l'un d'eux le pivot, et un pivot
sort de la liste. Ils restent ensemble jusqu'au premier pivot choisi parmi les rangs $i, i+1, \dots, j$. Si ce pivot est $i$ ou $j$, ils sont comparés. S'il se situe entre eux, ils sont séparés et ne sont jamais
comparés. Chaque nombre de ce bloc a la même probabilité d'être le premier, donc
\[
P(i \text{ et } j \text{ comparés}) = \frac{2}{j - i + 1}.
\]
En sommant sur toutes les paires, on obtient la formule. Pour $n = 10$, $2 \cdot 11 \cdot H_{10} - 40 \approx 24.44$.

\iqfollow{Comparez avec le pire cas et avec $n$ grand.}{Le pire cas vaut $n(n-1)/2 = 45$ pour $n = 10$, lorsque chaque pivot est le plus petit parmi ceux qui restent. Pour $n$ grand, la moyenne vaut environ $2n\ln n$, soit $1.39\, n \log_2 n$. Pour $n = 1000$, la moyenne exacte vaut environ 10\,986 comparaisons, contre 499\,500 dans le pire cas.}
\end{solution}

\begin{solution}{iq:iv1:expectation-and-indicators:22}
Le nombre moyen de donnes vaut exactement 10. Pour tout $n \ge 2$, il vaut exactement $n$.

L'idée clé est qu'une donne de $m$ cartes comporte en moyenne exactement une correspondance, quel que soit $m$, grâce aux variables indicatrices\index{variables indicatrices}. Ainsi, en moyenne, une carte sort à chaque donne, et $n$ donnes est la conjecture naturelle. Ce n'est pas encore une preuve, puisque le nombre de donnes est aléatoire, démontrons-le donc par
récurrence. Soit $p_k$ la probabilité d'avoir $k$ correspondances dans une donne de $n$ cartes, avec $\sum_k k\,p_k = 1$. Supposons $E[R_m] = m$ pour tous les $m$ strictement plus petits~; le cas $m = 1$ ne se présente jamais, puisqu'avoir exactement $n - 1$ correspondances est impossible. Par analyse au premier pas\index{analyse au premier pas},
\[
E[R_n] = 1 + p_0 E[R_n] + \sum_{k \ge 1} p_k (n - k) = 1 + p_0 E[R_n] + n(1 - p_0) - 1.
\]
D'où $(1 - p_0)E[R_n] = n(1 - p_0)$, et $E[R_n] = n$.

\iqwrong{Penser qu'il faut connaître les probabilités de dérangement. Ce n'est pas nécessaire~: elles s'annulent.}
\iqfollow{Combien de cartes individuelles distribuez-vous au total~?}{Chaque donne de $m$ cartes compte $m$ cartes distribuées. En reprenant la même analyse au premier pas avec $E[k] = 1$ et $E[k^2] = 2$, le total $S_n = n^2/2 + n$ vérifie la récurrence. Pour $n = 10$, cela fait 60 cartes.}
\end{solution}

\begin{solution}{iq:iv1:expectation-and-indicators:23}
L'espérance du nombre final vaut $e/2 - 1 \approx 0.359$.

Calculons en sommant selon la longueur de la suite croissante. Fixons $n \ge 1$ et une valeur $y$. La suite s'arrête au tirage $n+1$, celui-ci valant $y$, lorsque les $n$ premiers sont croissants et que le $n$-ième dépasse $y$. Les $n$ premiers sont croissants avec probabilité $1/n!$. Ils sont
croissants et tous inférieurs à $y$ avec probabilité $y^n/n!$. La probabilité recherchée vaut donc $(1 - y^n)/n!$, et
\[
E[\text{final}] = \sum_{n \ge 1} \int_0^1 y\,\frac{1 - y^n}{n!}\,dy = \sum_{n \ge 1} \frac{1}{n!}\Bigl(\frac12 - \frac{1}{n+2}\Bigr).
\]
La première partie vaut $(e-1)/2$. Pour la seconde, $\sum_{n \ge 0} \frac{1}{n!(n+2)} = \int_0^1 x e^x\,dx = 1$. En retirant le terme $n = 0$, qui vaut $\tfrac12$, il reste $\tfrac12$. D'où
\[
E[\text{final}] = \frac{e - 1}{2} - \frac12 = \frac{e}{2} - 1.
\]
\iqfollow{Et le nombre juste avant le dernier, le plus élevé de la suite~?}{La même méthode donne $\sum_{n \ge 1} \frac{1}{(n-1)!(n+2)} = \int_0^1 x^2 e^x\,dx = e - 2 \approx 0.718$.}
\iqfollow{Comment vérifiez-vous la réponse~?}{Le nombre de tirages a pour espérance $e$, comme pour toute suite croissante, et c'est un temps d'arrêt. L'identité de Wald\index{identité de Wald} indique donc que le total a pour espérance $e/2$. Une simulation confirme le total, le $0.359$ et le $0.718$.}
\end{solution}

\begin{solution}{iq:iv1:expectation-and-indicators:24}
L'espérance du nombre vaut $1 + \tfrac12 + \tfrac13 + \tfrac14 + \tfrac15 + \tfrac16 = 49/20 = 2.45$. En général, avec $n$ faces, c'est le nombre harmonique $H_n$.

Étiquetons les faces selon l'ordre de leur première apparition, et prenons la $k$-ième nouvelle face, $f$. Au moment où $f$ apparaît pour la première fois, $6 - k$ faces sont encore non vues. Alors $f$ n'apparaît qu'une seule fois si elle ne revient pas avant que toutes les faces non vues soient apparues.

Ignorons chaque lancer qui ne montre ni $f$ ni l'une de ces faces non vues\index{arguments de symétrie}. Les lancers restants sont uniformes sur $7 - k$ faces et indépendants du passé. Par symétrie, chacune de ces $7 - k$ faces a la même probabilité d'être la dernière à apparaître parmi les lancers futurs, donc $f$ est la dernière avec probabilité $1/(7-k)$. Si $f$ est la dernière, chaque face non vue apparaît avant sa prochaine apparition. Ainsi, la $k$-ième nouvelle face est un singleton avec probabilité $1/(7-k)$, et par linéarité\index{linéarité de l'espérance}
\[
E[\text{singletons}] = \sum_{k=1}^{6} \frac{1}{7-k} = H_6 = 2.45.
\]
\iqwrong{Utiliser un comptage de Poisson au temps d'arrêt espéré~: avec $14.7$ lancers, chaque face apparaît une fois avec une probabilité d'environ $(14.7/6)\,e^{-14.7/6}$, donc le total vaut environ $1.27$. Le temps d'arrêt dépend des comptages~: la règle arrête le jeu dès que la dernière nouvelle face apparaît, si bien que cette face est toujours un singleton. Traiter ce temps comme fixe passe à côté de cela.}
\iqfollow{Avec un jeu de 52 cartes, en tirant avec remise jusqu'à avoir vu toutes les cartes~?}{$H_{52} \approx 4.54$. Il croît très lentement, alors que l'attente totale croît comme $n \ln n$.}
\end{solution}

\begin{solution}{iq:iv1:expectation-and-indicators:25}
L'espérance du nombre de lancers vaut environ $24.13$.

La poissonisation\index{poissonisation} rend les faces indépendantes. Imaginons que les lancers arrivent selon un processus de Poisson de taux un par unité de temps. Chaque face apparaît alors comme un processus de Poisson indépendant de taux $1/6$. L'espérance du temps d'arrêt, en temps horloge, égale celle du nombre de lancers, car chaque lancer dure une unité en moyenne et la règle d'arrêt ne regarde que les lancers.

À l'instant $t$, avec $u = t/6$, une face donnée est apparue au moins deux fois avec probabilité $1 - e^{-u}(1 + u)$. Les six faces le sont toutes avec la puissance sixième de cette quantité. En intégrant la queue de la loi,
\[
E[\text{lancers}] = \int_0^\infty P(\tau > t)\,dt = 6 \int_0^\infty \Bigl[1 - \bigl(1 - e^{-u}(1 + u)\bigr)^6\Bigr] du.
\]
On développe par la formule du binôme. Chaque terme est une intégrale élémentaire~:
\[
E[\text{lancers}] = 6 \sum_{k=1}^{6} (-1)^{k+1} \binom{6}{k} \int_0^\infty e^{-ku}(1 + u)^k\,du,
\]
où l'intégrale pour $k = 1$ vaut $2$, celle pour $k = 2$ vaut $5/4$, et ainsi de suite. La somme vaut $24.13$.

Une alternative, que l'on peut traiter sous forme de tableau, est une chaîne de Markov sur $(a, b)$, où les deux indices comptent respectivement les faces vues zéro fois et une fois, avec $E(a,b) = \bigl[6 + a E(a-1, b+1) + b E(a, b-1)\bigr]/(a+b)$ et $E(0,0)=0$.

\iqwrong{En doublant l'attente de la collecte simple, on obtient $29.4$. Collecter le premier ensemble complet laisse des doublons, qui comptent pour le second. Il ne faut en moyenne que $9.43$ lancers de plus, au-delà des $14.7$.}
\end{solution}
